\chapter{\normalsize PENDAHULUAN}

%-----------------------------------------------------------------------------%
\section{\normalsize Latar Belakang}
%-----------------------------------------------------------------------------%
Kawasan strategis yang meliputi Provinsi Banten, DKI Jakarta, dan Jawa Barat tidak hanya berperan sebagai mesin pertumbuhan ekonomi nasional, tetapi juga menjadi arena dinamika pembangunan manusia yang sangat kompleks dalam lima tahun terakhir. Meskipun wilayah ini memberikan kontribusi yang sangat dominan terhadap perekonomian nasional dengan rata-rata pangsa Produk Domestik Regional Bruto (PDRB) berkisar 33,8 persen dari PDRB Indonesia \citet{bpspdrb}---keberhasilan pembangunannya tidak dapat lagi diukur hanya melalui kacamata indikator makro ekonomi semata. Terdapat interaksi multidimensi yang melibatkan kualitas hidup yang tecermin dalam Indeks Pembangunan Manusia (IPM), Indeks Pembangunan Gender (IPG), serta tantangan kesejahteraan sosial seperti kemiskinan dan ketahanan pangan yang terus mengalami perubahan. Kompleksitas ini menuntut pembuktian empiris yang lebih mendalam untuk melihat bagaimana indikator-indikator tersebut terdistribusi secara nyata di seluruh kabupaten/kota.

Secara empiris, lanskap pembangunan di ketiga provinsi ini menunjukkan polarisasi yang tajam antara pusat-pusat pertumbuhan utama dan wilayah penyangga. Meskipun data PDRB di kota-kota besar mencatat performa ekonomi yang kuat dengan nilai mencapai 532.606,87 miliar rupiah \citet{bpspdrb}, capaian tersebut tidak serta-merta diikuti oleh pemerataan kesejahteraan. Hal ini terlihat dari variasi spasial yang mencolok (\textit{striking spatial variation}) pada indikator kemiskinan, di mana Persentase Penduduk Miskin (PPM) terentang lebar dari 2,29 persen hingga mencapai 15,06 persen di wilayah lainnya. Kesenjangan ini juga terkonfirmasi oleh indikator daya beli, terlihat dari perbedaan Garis Kemiskinan (GK) dan Prevalensi Ketidakcukupan Konsumsi Pangan (PKKP) yang bervariasi lebar antarwilayah \cite{bpspkkp}. Lebih jauh lagi, disparitas kapabilitas manusia juga terlihat mencolok; Rata-rata Lama Sekolah (RLS) dan Angka Harapan Hidup (AHH) masih menunjukkan perbedaan yang nyata. Indeks Pembangunan Gender (IPG) serta Tingkat Pengangguran Terbuka (TPT) masih berada di kisaran 6,55--10,12 persen pada rentang tahun 2020--2024. Namun, persoalan ketimpangan antarwilayah ini menjadi semakin rumit ketika diletakkan dalam konteks waktu yang penuh perubahan.

Selain disparitas spasial, tantangan analisis meningkat akibat adanya dinamika temporal yang sangat fluktuatif sepanjang periode 2020--2024. Periode ini diawali dengan peristiwa pandemi COVID-19 yang menekan indikator makroekonomi dan kesejahteraan secara serentak terlihat dari perubahan pertumbuhan ekonomi yang sempat mencapai angka 1,88 persen pada tahun 2020 \citet{bpspdrb} namun diikuti oleh pola pemulihan yang lajunya tidak seragam di setiap daerah pada tahun-tahun berikutnya. Perubahan drastis dari masa pandemi COVID-19 menuju masa pemulihan ini diduga kuat telah mengubah struktur hubungan antar-indikator secara fundamental; sebagai contoh, pola korelasi antara pertumbuhan ekonomi (PDRB) dengan tingkat pengangguran (TPT) saat krisis sangat mungkin berbeda dibandingkan saat kondisi normal kembali pulih pada tahun 2023--2024. Oleh karena itu, pendekatan analisis multivariat klasik yang mengasumsikan struktur hubungan variabel bersifat statis (tetap sepanjang waktu) menjadi tidak lagi relevan untuk menganalisis fenomena pada periode yang penuh fluktuasi ini.

Untuk mengatasi keterbatasan asumsi stabilitas tersebut, diperlukan pendekatan analisis data tiga-arah (\textit{three-way data analysis}) yang mampu menangani struktur data kompleks yang terdiri dari objek (kabupaten/kota), variabel (indikator sosio-ekonomi), dan situasi (waktu). Penelitian ini mengusulkan penggunaan metode STATIS (\textit{Structuration des Tableaux à Trois Indices de la Statistique}), sebuah metode eksploratif yang diperkenalkan oleh \citet{lavit94,oliveira2024}. Berbeda dengan Analisis Komponen Utama (PCA) konvensional yang hanya bekerja pada data dua dimensi, STATIS dirancang khusus untuk mengekstraksi informasi dari sekumpulan tabel data yang berubah seiring waktu. Metode ini memiliki keunggulan dalam membentuk sebuah struktur kompromi (\textit{compromise structure}) yang merepresentasikan pola hubungan umum antar-indikator selama periode pengamatan, sekaligus memungkinkan identifikasi seberapa jauh struktur ekonomi setiap tahunnya menyimpang dari pola umum tersebut.

Meskipun metode STATIS mampu menghasilkan struktur kompromi secara matematis, kompleksitas hasil analisis dari matriks berdimensi tinggi sering kali sulit dipahami jika hanya disajikan dalam bentuk tabel numerik. Oleh karena itu, penelitian ini mengintegrasikan pendekatan STATIS dengan teknik visualisasi \textit{Dynamic Biplot}. Pendekatan ini memungkinkan proyeksi posisi objek (kabupaten/kota) dan variabel (indikator ekonomi) ke dalam ruang berdimensi rendah, sekaligus menggambarkan pergerakan atau lintasan (\textit{trajectory}) perubahan posisi tersebut dari tahun ke tahun. Dengan memanfaatkan konsep \textit{trajectory}, pergeseran karakteristik suatu wilayah misalnya, sebuah kabupaten yang pada tahun 2020 memiliki tingkat pengangguran tinggi namun kemudian bergeser mendekati karakteristik wilayah dengan IPM tinggi pada tahun 2024 dapat dilacak secara visual. Visualisasi lintasan ini menjadi instrumen krusial untuk membedakan antara fluktuasi jangka pendek akibat pandemi dan perubahan struktural jangka panjang.

Penerapan metode eksploratif multivariat seperti ini mengisi celah penting dalam literatur analisis regional di Indonesia yang selama ini didominasi oleh pendekatan inferensial klasik. Studi-studi terdahulu umumnya berfokus pada pemodelan regresi data panel atau model spasial ekonometrik untuk mengestimasi besaran pengaruh variabel independen terhadap variabel dependen. Sebagai contoh, penelitian yang dilakukan oleh \citet{nuswantari2021} menggunakan pendekatan \textit{Dynamic Spatial Durbin Model} (SDM) untuk memodelkan kemiskinan, sedangkan \citet{ahmad2024} menerapkan pendekatan spasial yang serupa untuk menganalisis fenomena penyebaran penyakit kusta di Pulau Jawa. Fokus utama dari penelitian-penelitian tersebut adalah pengukuran besaran pengaruh peubah-peubah penjelas (\textit{explanatory variables}) terhadap peubah respon yang terbukti nyata secara statistik. Meskipun pendekatan tersebut berhasil mengidentifikasi faktor-faktor determinan, metode inferensial ini memiliki keterbatasan dalam memvisualisasikan struktur hubungan antarvariabel secara simultan (\textit{inter-variable structure}) dan bagaimana struktur tersebut berevolusi dari waktu ke waktu (\textit{spatio-temporal}). Hingga saat ini, belum terdapat penelitian yang secara spesifik menggunakan pendekatan eksploratif untuk memetakan dinamika struktur multivariat indikator ekonomi secara komprehensif tanpa membedakannya menjadi variabel bebas dan terikat.

Kompleksitas informasi yang dihasilkan dari analisis spasio-temporal ini menuntut media penyajian yang lebih adaptif dibandingkan laporan statis konvensional. Untuk itu, penelitian ini tidak hanya berhenti pada tataran analisis teoritis, tetapi juga mengembangkan sebuah aplikasi visualisasi interaktif berbasis web menggunakan bahasa pemrograman Python dengan kerangka kerja (\textit{framework}) Dash dan Plotly. Pengembangan \textit{dashboard} ini bertujuan untuk memfasilitasi Eksplorasi Data (\textit{Exploratory Data Analysis}) secara mandiri oleh pengguna. Fitur-fitur interaktif yang disematkan seperti seleksi tahun, filter wilayah, hingga rotasi \textit{baseline} memungkinkan para pengambil kebijakan untuk mengamati pola interaksi multivariat secara lebih mendalam. Melalui antarmuka ini, pola-pola tersembunyi (\textit{latent patterns}) yang mungkin terlewatkan dalam tabel statistik standar dapat terdeteksi dengan lebih intuitif.

Secara keseluruhan, penelitian ini diarahkan untuk menjawab kebutuhan akan metode pemantauan pembangunan daerah yang lebih menyeluruh dan dinamis. Dengan memadukan kekuatan analitis metode STATIS dalam memahami struktur kompromi (\textit{compromise structure}) lintas tahun dan kemampuan naratif visual dari \textit{Dynamic Biplot}. Penelitian ini diharapkan dapat memberikan kontribusi ganda: secara akademis memperkaya khazanah penerapan statistik multivariat pada data ekonomi regional, dan secara praktis menyediakan purwarupa alat bantu visual yang relevan bagi pemerintah daerah dalam mengevaluasi dampak kebijakan pasca-pandemi. Identifikasi struktur hubungan yang stabil maupun yang berubah selama lima tahun terakhir akan menjadi landasan penting dalam merumuskan strategi pembangunan yang lebih tangguh di masa depan.

Pemilihan lokasi penelitian yang mencakup seluruh kabupaten/kota di Provinsi Banten, DKI Jakarta, dan Jawa Barat didasarkan pada karakteristik interkonektivitas ekonomi yang kuat di antara ketiga wilayah ini. Sebagai kesatuan wilayah pembangunan utama di Pulau Jawa, ketiga provinsi ini memiliki keterkaitan rantai pasok, mobilitas tenaga kerja, dan arus investasi yang sangat intensif. Permasalahan ekonomi yang terjadi di pusat pertumbuhan utama (DKI Jakarta) diproyeksikan akan merambat (\textit{spillover}) ke wilayah penyangga di Banten dan Jawa Barat dengan skala yang berbeda-beda. Oleh karena itu, analisis yang hanya membatasi lingkup pada wilayah metropolitan Jabodetabek berpotensi mengabaikan dampak sistemik yang dirasakan oleh wilayah kabupaten non-metropolitan yang secara administratif berada dalam satu kesatuan provinsi namun memiliki karakteristik ekonomi yang jauh berbeda.

Untuk memahami kompleksitas pembangunan tersebut, penelitian ini menggunakan dua belas indikator strategis yang merepresentasikan dimensi pembangunan secara menyeluruh. Pemilihan variabel ini tidak dilakukan tanpa alasan, melainkan dipertimbangkan untuk mencakup aspek makroekonomi (PDRB ADHK), kesejahteraan material (PPM, GK, PKKP), kualitas modal manusia (IPM, RLS, AHH Laki-laki dan Perempuan), kesetaraan (IPG, Gini Ratio), serta partisipasi pasar kerja (TPT, TPAK). Kombinasi variabel ini memungkinkan analisis untuk tidak hanya pada indikator pertumbuhan konvensional, namun juga aspek inklusivitas pembangunan. Sebagai contoh, analisis simultan antara Gini Ratio dan PDRB memungkinkan deteksi dini terhadap fenomena pertumbuhan yang tidak merata (\textit{exclusive growth}), sementara interaksi antara TPAK dan IPG dapat mengungkap hambatan struktural dalam partisipasi angkatan kerja berbasis gender.

Ketiadaan studi yang menganalisis dinamika interaksi kedua belas indikator ini secara spasio-temporal pada periode kritis 2020--2024 menciptakan sebuah "kotak hitam" (\textit{black box}) dalam pemahaman mengenai resiliensi ekonomi regional. Tanpa analisis struktur kompromi dan visualisasi lintasan yang jelas, para pemangku kebijakan kehilangan panduan empiris mengenai bagaimana sebenarnya pola pemulihan berlangsung di tingkat lokal apakah kembali ke struktur lama atau membentuk ekuilibrium baru. Berdasarkan urgensi untuk mengungkap fenomena tersebut secara statistik dan visual, penelitian ini disusun untuk menjawab pertanyaan-pertanyaan mendasar mengenai evolusi struktur ekonomi dan pembangunan di ketiga provinsi tersebut.

%-----------------------------------------------------------------------------%
\section{\normalsize Rumusan Masalah}
%-----------------------------------------------------------------------------%
Berdasarkan latar belakang yang telah diuraikan, permasalahan yang akan dikaji dalam penelitian ini dirumuskan sebagai berikut:

\begin{enumerate}
    \item Bagaimana karakteristik struktur kompromi (\textit{compromise structure}) yang merepresentasikan pola hubungan multivariat dua belas indikator pembangunan di seluruh kabupaten/kota Provinsi Banten, DKI Jakarta, dan Jawa Barat menggunakan metode STATIS untuk memahami pola umum lintas tahun?
    \item Bagaimana menguji struktur hubungan antar-indikator tersebut secara statistik selama periode 2020--2024 menggunakan uji permutasi matriks inter-struktur untuk membuktikan keberadaan dinamika temporal?
    \item Bagaimana pola lintasan (\textit{trajectory}) perubahan karakteristik ekonomi kabupaten/kota dari masa krisis pandemi (2020) hingga masa pemulihan (2024) berdasarkan visualisasi \textit{Dynamic Biplot}?
    \item Bagaimana merancang aplikasi berbasis web menggunakan \textit{framework} Dash dan Plotly sebagai alat bantu eksplorasi data yang memungkinkan pengguna melakukan pemantauan indikator pembangunan secara interaktif?
\end{enumerate}

%-----------------------------------------------------------------------------%
\section{\normalsize Ruang Lingkup Penelitian}
%-----------------------------------------------------------------------------%
Agar pembahasan penelitian ini lebih terarah dan fokus, penulis menetapkan batasan-batasan masalah sebagai berikut:

\begin{enumerate}
    \item \textbf{Variabel Penelitian:} Analisis menggunakan 12 indikator terpilih yang bersumber dari Badan Pusat Statistik (BPS), meliputi: Indeks Pembangunan Manusia (IPM), Indeks Pembangunan Gender (IPG), Persentase Penduduk Miskin (PPM), PDRB ADHK, Rata-rata Lama Sekolah (RLS), Angka Harapan Hidup Laki-laki (AHH\_L) dan Perempuan (AHH\_P), Garis Kemiskinan (GK), Prevalensi Ketidakcukupan Konsumsi Pangan (PKKP), Tingkat Pengangguran Terbuka (TPT), Tingkat Partisipasi Angkatan Kerja (TPAK), dan Gini Ratio (GR).
    \item \textbf{Aspek Spasial:} Istilah "spasial" dalam penelitian ini merujuk pada unit analisis berbasis wilayah (kabupaten/kota), bukan pada penggunaan matriks pembobot spasial (\textit{spatial weight matrix}) seperti pada Ekonometrika Spasial.
    \item \textbf{Fungsionalitas Aplikasi:} Aplikasi berbasis web yang dikembangkan berfungsi sebagai alat bantu eksplorasi visual (\textit{Visual Exploratory Tool}) bagi pengguna, bukan sebagai Sistem Pendukung Keputusan (\textit{Decision Support System}) yang menghasilkan rekomendasi kebijakan berbasis data secara otomatis.
\end{enumerate}

%-----------------------------------------------------------------------------%
\section{\normalsize Tujuan Penelitian}
%-----------------------------------------------------------------------------%
Sejalan dengan rumusan masalah di atas, tujuan dari penelitian ini adalah:

\begin{enumerate}
    \item Menganalisis karakteristik struktur kompromi (\textit{compromise structure}) yang merepresentasikan pola hubungan multivariat dua belas indikator pembangunan di seluruh kabupaten/kota Provinsi Banten, DKI Jakarta, dan Jawa Barat menggunakan metode STATIS untuk memahami pola umum lintas tahun.
    \item Menganalisis struktur hubungan antar-indikator tersebut secara statistik selama periode 2020--2024 menggunakan uji permutasi matriks inter-struktur untuk membuktikan keberadaan dinamika temporal.
    \item Menganalisis pola lintasan (\textit{trajectory}) perubahan karakteristik ekonomi kabupaten/kota dari masa krisis pandemi (2020) hingga masa pemulihan (2024) berdasarkan hasil visualisasi \textit{Dynamic Biplot}.
    \item Merancang aplikasi berbasis web menggunakan \textit{framework} Dash dan Plotly sebagai alat bantu eksplorasi data yang memungkinkan pengguna melakukan pemantauan indikator pembangunan secara interaktif.
\end{enumerate}

%-----------------------------------------------------------------------------%
\section{\normalsize Metodologi Penelitian}
%-----------------------------------------------------------------------------%
Penelitian ini menggunakan pendekatan eksploratif multivariat (\textit{Descriptive Multivariate Analysis}). Penelitian tidak bertujuan untuk membangun model prediktif (peramalan masa depan) maupun model inferensial kausalitas (seperti uji pengaruh parsial dalam regresi).